\documentclass[10pt]{article}
\usepackage[hmargin=20mm,vmargin=20mm]{geometry}
\usepackage{amsmath, amsthm, amssymb, mathabx, verbatim, setspace, enumerate}
\usepackage{relsize}
\usepackage[mathscr]{euscript}
\usepackage[all,cmtip]{xy}
\usepackage{relsize}
\usepackage[toc,page]{appendix}

\usepackage{tikz-cd}

\xyoption{rotate}
\xyoption{pdf}


\usepackage{tikz-cd}


\linespread{1.1}

%prevents underfill for URLs in references
\usepackage{url}
\Urlmuskip=0mu plus 1mu\relax

%hyperlinked ToC
\usepackage{hyperref}



\newtheorem{thm}{Theorem}[section]
\newtheorem{prop}[thm]{Proposition}
\newtheorem{lemma}[thm]{Lemma}
\newtheorem{cor}[thm]{Corollary}

\theoremstyle{definition}
\newtheorem{exmp}[thm]{Example}
\newtheorem{defn}[thm]{Definition}
\newtheorem{rmk}[thm]{Remark}
\newtheorem{notation}[thm]{Notation}
\newtheorem{exercise}{Exercise}
\newtheorem{lem}[thm]{Lemma}

\DeclareMathOperator{\Stab}{Stab}
\DeclareMathOperator{\Hom}{Hom}
\newcommand{\colim}{\operatornamewithlimits{colim}}
\DeclareMathOperator{\Sh}{Sh}
\DeclareMathOperator{\Pro}{Pro}
\DeclareMathOperator{\Ind}{Ind}
\DeclareMathOperator{\Ext}{Ext}
\DeclareMathOperator{\ch}{ch}
\DeclareMathOperator{\supp}{supp}
\DeclareMathOperator{\Coh}{Coh}
\DeclareMathOperator{\dmod}{-mod}
\DeclareMathOperator{\dperf}{-perf}
\DeclareMathOperator{\Ho}{Ho}
\DeclareMathOperator{\dcmod}{-cmod}
\DeclareMathOperator{\Spec}{Spec}
\DeclareMathOperator{\im}{im}
\DeclareMathOperator{\gr}{gr}
\DeclareMathOperator{\Sym}{Sym}
\DeclareMathOperator{\shHom}{\mathcal{H}\emph{om}}
\DeclareMathOperator{\shExt}{\mathcal{E}\emph{xt}}
\DeclareMathOperator{\Mod}{Mod}
\DeclareMathOperator{\DMod}{DMod}
\DeclareMathOperator{\pt}{pt}
\DeclareMathOperator{\defect}{def}
\DeclareMathOperator{\codim}{codim}
\DeclareMathOperator{\coker}{coker}
\DeclareMathOperator{\dR}{dR}
\DeclareMathOperator{\grr}{gr}
\DeclareMathOperator{\opp}{op}
\DeclareMathOperator{\Aff}{Aff}
\DeclareMathOperator{\Tor}{Tor}
\DeclareMathOperator{\Der}{Der}
\DeclareMathOperator{\End}{End}
\DeclareMathOperator{\Res}{Res}
\DeclareMathOperator{\QCoh}{QCoh}
\DeclareMathOperator{\tr}{tr}
\DeclareMathOperator{\DCoh}{DCoh}
\DeclareMathOperator{\Perf}{Perf}
\DeclareMathOperator{\Frac}{Frac}
\DeclareMathOperator{\IndCoh}{IndCoh}
\DeclareMathOperator{\Proj}{Proj}
\DeclareMathOperator{\Rep}{Rep}
\DeclareMathOperator{\cone}{cone}
\DeclareMathOperator{\Tot}{Tot}
\DeclareMathOperator{\ICoh}{ICoh}
\DeclareMathOperator{\Tate}{Tate}
\DeclareMathOperator{\Map}{Map}
\DeclareMathOperator{\Fun}{Fun}
\DeclareMathOperator{\LMod}{LMod}


%Complicated commands
\newcommand{\leftexp}[2]{{\vphantom{#2}}^{#1}{#2}}

%Sheafy commands
\newcommand{\sEnd}{\mathcal{E}\emph{nd}}


%Shorthand commands
\newcommand{\dd}{\partial}
\newcommand{\ra}{\rightarrow}
\newcommand{\la}{\leftarrow}
\newcommand{\mf}{\mathfrak}
\newcommand{\bs}{\backslash}
\newcommand{\cat}{\mathbf}
\newcommand{\wh}{\widehat}
%mathbb
\newcommand{\PP}{\mathbb{P}}
\newcommand{\Z}{\mathbb{Z}}
\newcommand{\C}{\mathbb{C}}
\newcommand{\R}{\mathbb{R}}
\newcommand{\N}{\mathbb{N}}
\newcommand{\A}{\mathbb{A}}
\newcommand{\G}{\mathbb{G}}
\newcommand{\Q}{\mathbb{Q}}
%mathcal
\newcommand{\OO}{\mathcal{O}}
\newcommand{\DD}{\mathcal{D}}
\newcommand{\LL}{\mathcal{L}}
\newcommand{\LLf}{\widehat{\LL}}
%mathfrak
\newcommand{\h}{\mf{h}}
\newcommand{\g}{\mf{g}}



\DeclareMathOperator{\Aut}{Aut}
%Typo correction
\newcommand{\tidle}{\tilde}
\newcommand{\rightarrw}{\rightarrow}
\newcommand{\alhpa}{\alpha}
\newcommand{\rigtarrow}{\rightarrow}
\newcommand{\emoh}{\emph}
\newcommand{\codts}{\cdots}
\newcommand{\shEnd}{\sEnd}
\newcommand{\rihgtarrow}{\rightarrow}


%Markup
\newcommand{\todo}[1]{\textbf{\textcolor{red}{TODO #1}}}

\newcommand{\loccit}{\emph{loc. cit. }} 

\setcounter{tocdepth}{2}


\title{Homework 2, due 9/20}
\author{Math 6670}

\begin{document}
\maketitle

\begin{enumerate}
\item Prove that the homotopy category of a pretriangulated dg category over a ring $k$ is naturally a triangulated category by defining the shift functor $T(X) = \cone(X \rightarrow 0)$ and by defining the class of exact triangles as follows.  Let $\Delta^1$ denote the dg category consisting of two objects: $0$ and $1$, such that $\Hom(0, 0) \simeq \Hom(1, 1) \simeq k \cdot \{\mathrm{id}\}$ and $\Hom(0, 1) = k \cdot \{f\}$.  Let $\cat{C}^\Delta_{uni}$ denote the full subcategory in its Yoneda embedding containing $h_0, h_1, h_{\cone(f)}$.
\begin{enumerate}
\item Show that $\Fun(\Ho(\cat{C}^\Delta_{uni}), \Ho(\cat{C}))$ is the category of triangles in $\Ho(\cat{C})$, and that $\Fun_{dg}(\cat{C}^\Delta_{uni}, \cat{C})$ is a subcategory (is it full?).
\item Let $\cat{C}^\Delta_{rot}$ be the rotationally symmetric model for $\cat{C}^\Delta_{uni}$ discussed in class, i.e. with objects $X, Y, Z$, morphisms $f: X \rightarrow Y$, $g: Y \rightarrow Z$, and $h: Z \rightarrow X[1]$, $F: Y \rightarrow X$, $G: Z \rightarrow Y$, $H: X[1] \rightarrow Z$, such that $d(f) = d(g) = d(h) = 0$ and $d(F) = hg$, $d(G) = fh$, $d(H) = gf$, and $Ff - gG = \mathrm{id}_Y$, $Gg - hH = \mathrm{id}_Z$, and $Hh - fF = \mathrm{id}_X$.  This category evidently has a $C_3$-action.  Construct a dg functor $\cat{C}^\Delta_{uni} \rightarrow \cat{C}^\Delta_{rot}$ which is an equivalence of dg categories.
\item Define the class of exact triangles to be those triangles which come from a dg functor.  Show that this class satisfies the axioms of a triangulated category, except for the octahedral axiom.
\end{enumerate}


\item Huybrechts, Exercise 2.20.  Show that a complex $A^\bullet$ is isomorphic to $0$ in $\mathcal{D}(\mathcal{A})$ if and only if $H^i(A^\bullet) \simeq 0$ for all $i$.  Find an example of a complex morphism $f: A^\bullet \rightarrow B^\bullet$ such that $H^i(f) = 0$ for all $i$, but $f$ is not trivial in $\mathcal{D}(\mathcal{A})$.
\item Let $\cat{A}$ be a category with enough injectives, and let $C_{inj}(\cat{A})$ denote the full subcategory of $C(\cat{A})$ of complexes of injective objects and $K_{inj}(\cat{A})$ its homotopy category.  Show that $C_{inj}(\cat{A})$ is pretriangulated, so that $K_{inj}(\cat{A})$ is triangulated.
\begin{enumerate}
\item Show that if $A^\bullet$ is a bounded below (left bounded) complex such that $H^i(A^\bullet) = 0$ for all $i$, then $A^\bullet \simeq 0$ in $K_{inj}(\cat{A})$.
\item Find an example of an injective complex $A^\bullet$ which is has $H^i(A^\bullet) = 0$ for all $i$, but $A^\bullet$ is not equivalent to 0 in $K_{inj}(\cat{A})$.  Hint: Take $\cat{A} = R\dmod$ for $R$ a Frobenius algebra so that free modules are injective.
\end{enumerate}
\item Let $\cat{A}$ be a hereditary abelian category, i.e. $\Ext^2$ vanishes for all objects.  Let $H(\cat{A}) \subset C(\cat{A})$ denote the full subcategory of complexes whose internal differentials vanish (i.e. formal complexes).  Define a dg functor $H(\cat{A}) \rightarrow D(\cat{A})$ which is an equivalence of dg categories and prove it.
\item Let $F: \cat{A} \rightarrow \cat{B}$ be a functor of abelian categories which is left-exact.  We say that a class of objects $\cat{K}_F \subset \cat{A}$ is \emph{$F$-adapted} if it is closed under finite direct sums, if $F$ sends acyclic objects of $K^+(\cat{A})$ whose terms are in $\cat{K}_F$ to acyclic objects and if any object of $\cat{A}$ can be embedded into an object of $\cat{K}_F$. 
\begin{enumerate}
\item Show, using the universal definition of exact functors (e.g. in Haiman's notes), that $RF(A^\bullet)$ is well-defined and is computed by an object in $\cat{K}_F$.
\item Let $X$ be a Noetherian and separated scheme over a ring $R$, and let $\Gamma: \QCoh(X) \rightarrow R\dmod$ be the global sections functor.  Find a $\Gamma$-adapted class of objects such that the terms of the Cech complex are $\Gamma$-adapted.
\item Conclude that the Cech complex computes $R\Gamma(X, -)$.
\end{enumerate}
\end{enumerate}

\end{document}
